\chapter{Learning architectures}\label{ch:architecture}
Blackbox methods for lifelong learning are promising approaches to achieve continually adaptive control, but are hard-to-interpret.\ What other techniques can be implemented that are more intuitive in nature?

Well-designed control and learning architectures can either enable or facilitate lifelong learning, whether the control schemes use blackbox models or more traditional methods.\ For example, some techniques use two or more blackbox models, where the idea is to have each model deal with a separate challenge of the lifelong learning goal.\ Other architectural methods aim to expand existing control schemes or policies in order to enable them to continually adapt.

In this chapter, these two main strategies will be further explored in the literature. % TODO: update sentence (sec on trad methods)
\comment{Added introduction to chapter}


\section{Traditional adaptive control architectures}\label{sec:adaptive-control}



\section{Dual architectures}\label{sec:dual}
Under `dual architectures', the methods that have separate modules for solving different aspects of the lifelong learning challenge are gathered.\ The term `dual' is used here in a very broad sense because each module might ofcourse be a grouping of various techniques with intricate architectures.\ However, studying the literature, it can be found that most of the methods discussed here can be broken down into two main modules.\ Generally, these dual architectures often exist either of a policy/control module in combination with some type of knowledge base used for storing previous experiences and informing the policy (refer to \autoref{ch:storage} for details on how knowledge bases can differ in implementation), or two (or more) control policies that are each in turn capable of handling low uncertainty and high uncertainty situations.

To illustrate how these modules can be a very broad concept, consider a knowledge base: A knowledge base must first and foremost store knowledge, but additionally also maintain and update this library.\ Furthermore, the knowledge base should be capable of categorizing new observations in relation to its library.\ All these capabilities are seldom performed by a single module, but the techniques used vary wildly between different implementations.\ It is therefore not very interesting to differentiate between each different implementation in this chapter, but rather use the overarching term `knowledge base' and discuss the more intricate details in \autoref{ch:storage}.

Having established the broad framework and significance of dual architectures, examples from literature that illustrate how these modular approaches have been effectively implemented can be discussed.

In \textcite{Wang2022}, a knowledge base stores potentially infinite sets of parametrized environments along with their respective control strategies.\ This repository is used to guide the current control policy: When a new environment is encountered, the environment models in the knowledge base are evaluated on how well they fit the observations.\ The control policy is then initialized using the policy parameters belonging the most similar environment found in this knowledge base and the policy continues to learn online.\ In addition, environment models can be added to the library when no suitable match to the observations can be found.

\textcite{Schwarz2018} utilize an architecture consisting of two modules, which they name the knowledge base and the active column.\ During two alternating phases, progression and compression, the active column learns a new task and the learned policy is distilled into the knowledge base, respectively.\ The knowledge base in this case is a policy $\pi_{t}^{KB}$ that serves as a stable reference for the active policy $\pi_{t+1}$ during progression and remains unchanged at this time.\ During compression the newly learned active policy is distilled into $\pi_{t}^{KB}$ resulting in an updated $\pi_{t+1}^{KB}$.

In order to combat the issue of forgetting previous policies while training on new tasks, a common approach is to store old experiences and intermittently reuse them during the training process of a new policy.\ One downside is that this approach can be very memory-inefficient.\ \textcite{Shin2017} therefore propose a method where a deep generative model, a so-called \emph{Generator}, is trained to mimic the training data provided for a task-solving model called the \emph{Solver}.\ The Solver is then trained on a mix of actual training data and generated samples provided by the Generator.\ Likewise, when a new Generator is trained to mimic the data of the next task, some samples created by the previous Generator are mixed in.\ This ensures that knowledge is not forgotten, since each Generator also learns to recreate data generated by its predecessor.

An example of another approach generally taken when using dual architectures is found in \textcite{Lu2019}.\, Two modules, a model-free and model-based module, work in synthesis to simulate the way human beings behave: functioning on \emph{autopilot} most of the time and extensively planning only when we are uncertain.\ In the model-free module, a policy optimisation algorithm (in this case TD3\todo{citation [14] from \textcite{Lu2019}}) rolls out an initial plan, while the standard deviation over an ensemble of value functions is used to provide measure of uncertainty.\ Based on this uncertainty a planning horizon $H$ is set and the model-based planner (MPPI) now starts updating the initial plan for the given horizon, using the ensemble of value functions to evaluate its plans.\ During this model-based planning, any rollouts of potential plans are added to a replay buffer $\mathcal{D}_{\pi}$ which is in turn used to train the model-free policy optimization algorithm.

\todo{\cite{Lu2020}} % TODO: Reset-free LL w/ skill-space planning

\todo{\cite{Zhou2022}} % TODO: Forgetting and imbalance in Robot LL with off-policy data


\section{Architectural expansions}\label{sec:expansions}
`Architectural expansions' is a grouping of methods where an already functional control scheme or policy is enhanced by adding one or more modules to the scheme, enabling it to learn and/or adapt continually without actually altering the preexisting policy.\ The idea behind such techniques was already discussed in \autoref{sec:adaptive-control}, but here the focus will be on approaches that incorporate more of a learning aspect rather than simply adapting.
~\cite{Chatzilygeroudis2016, Nazeer2023, Sun2011}.

\todo{\cite{Sprechmann2018}} % TODO: Memory-based parameter adaptation